\chapter{Hypothesis and Objectives}
\chapterquote{Always look forward, keep your eyes on the goal, and do not stop until you reach it!}{All Might, My Hero Academia}


\newpage
\section{Hypothesis}

Because the molecular mechanisms underlying AMP function occur at complex interfaces, their transient and dynamic nature makes the precise characterization of these membrane-disrupting processes a significant experimental challenge\cite{hypothesis,srm_methods}. In this context, unbiased molecular dynamics (MD) and enhanced sampling techniques\cite{opes,steered_md,umbrella_sampling}, such as metadynamics\cite{Barducci2008}, provide a powerful complementary approach\cite{hypothesis}, offering precise atomistic and thermodynamic insights into how peptide orientation, specifically tilt and spin angles, drives selective membrane recognition and lipid assembly integration.  

A second hypothesis underlying this work is that the combinatorial explosion inherent to interfacial peptide folding and environment-conditioned sequence design; a problem characterized by a vast, high-dimensional space of potential solutions\cite{Levinthal_paradox,Levinthal_paradox_2}; can be effectively addressed through the application of Quantum Computing (QC)\cite{biology_quantum_computing,robert2021resource,perdomo_1,perdomo_2}. Classical computational approaches often struggle with the exponential scaling required to exhaustively explore these immense sequence and conformational landscapes\cite{biology_quantum_computing}. To overcome this fundamental bottleneck, this thesis posits that formulating the problem via physics-based Hamiltonian models allows the complex energetic landscape to be mapped into compact binary encodings. This quantum-ready framework enables the deployment of emerging hybrid quantum-classical algorithms, such as the Variational Quantum Eigensolver (VQE)\cite{peruzzo2014variational} and the Quantum Approximate Optimization Algorithm (QAOA)\cite{farhi2014quantum}. By leveraging these quantum techniques, it is possible to navigate the astronomical number of potential solutions more efficiently than classical exact enumeration, thereby accelerating the rational design of novel AMPs with tailored environmental specificity.


\section{General and specific Objectives}
Driven by these hypotheses, the overarching objective of this doctoral thesis is to decipher the molecular determinants that govern the sequence-structure-function relationships of $\alpha$-helical antimicrobial peptides (AMPs). Ultimately, this research aims to predict their interfacial folding dynamics and selective affinity for pathological lipid membranes, specifically those implicated in the Cancer-Aging-Infection (CAIn) triangle. To systematically achieve this goal, the work was structured around the following specific objectives:

\begin{enumerate}[font=\bfseries]
    \item \textbf{Systematic biophysical characterization of AMP-membrane interactions.} Previous studies often lack a detailed thermodynamic and mechanistic description of how AMPs insert into varying lipid environments[cite: 3]. The first objective of this thesis was to employ unbiased molecular dynamics (MD) and well-tempered metadynamics to decode the interaction of canonical AMPs (e.g., LL-37, magainin-2, pleurocidin) with distinct lipid bilayers mimicking healthy mammalian, bacterial, and cancerous cells[cite: 3]. This involved calculating accurate Free Energy Surfaces (FES) and determining the decisive role of the transverse hydrophobic dipole moment, tilt, and spin angles in governing peptide insertion and orientation[cite: 3]. This objective establishes the biophysical foundation of the thesis and is addressed in the first publication presented in chapter \ref{chap1}.
    
  
    
    \item \textbf{Implementation of peptide folding and design algorithms on Quantum Computers.} To overcome the exponential scaling bottlenecks of classical exhaustive enumeration for sequence design and structural prediction, the third objective focused on translating the developed physical models into quantum-compatible formats[cite: 1, 4]. This involved mapping the environment-conditioned Hamiltonian into a compact binary encoding and utilizing a tetrahedral lattice model[cite: 1, 4]. Subsequently, hybrid quantum-classical algorithms, such as the Variational Quantum Eigensolver (VQE) and the Quantum Approximate Optimization Algorithm (QAOA), were successfully applied to predict the optimal folding conformations of peptides explicitly positioned at the smooth interface between polar and non-polar media, and to optimize sequences[cite: 1, 4]. This objective pushes the boundaries of current computational biophysics by embracing emerging technologies, and is presented in the third publication in chapter \ref{chap2}.
    
      \item \textbf{Formulation of a physics-based Hamiltonian for environment-conditioned sequence design.} Designing peptide sequences that remain stable and selective across heterogeneous environments is a major computational challenge[cite: 1]. The second objective of this thesis addressed this limitation by developing a rigorous energy function that integrates helix propensities, pairwise interactions, electrostatics, solvent polarity, and interfacial geometry[cite: 1]. By normalizing these physical contributions as Z-scores relative to random sequence ensembles, this framework enables the quantitative evaluation and rational design of AMP sequences with tunable specificity for aqueous, apolar, or anionic membrane interfaces[cite: 1]. This objective drives the transition from observation to rational design, and is addressed in the second publication of this thesis, included in chapter \ref{chap3}.
    
    \item \textbf{Development of accessible computational tools for the peptide community.} The final objective of this thesis was to democratize access to the advanced computational methodologies generated throughout this research[cite: 1]. To this end, web-based tools were developed to allow non-expert users to access environment-conditioned design frameworks and sequence-landscape diagnostics in a straightforward and user-friendly manner[cite: 1]. This objective emphasized the translation of methodological advances into practical research tools, culminating in the deployment of the HelixDesignStudio web implementation[cite: 1]. The fulfillment of this objective ensures the broader impact of this work and is detailed in the final publication included in chapter \ref{chap3}.
\end{enumerate}

The designed tasks were expected to provide a deep understanding of how amino acid sequence dictates AMP selectivity towards pathological membranes. Furthermore, to facilitate access to sequence design, computational scoring, and hybrid quantum-classical workflows, the tools generated during this thesis were integrated into accessible web platforms, aiming to accelerate the discovery of novel targeted therapies within the context of the CAIn triangle.